\section{Local selection rules and energy scales}
\label{sec:local-details}

Figure~\ref{fig:selection} collects the local projection and selection-rule
data for positive $\alpha$.
At $\alpha=0.14,0.23,0.33$, the endpoint projections give
$\Jeff/\beta=0.440,0.432,0.419$, respectively.
The lowest odd-parity excited doublet carries vanishing $q=\pi$ weight.
Weights are summed over complete final multiplets and averaged over the
initial ground multiplet; the $84\%$--$98\%$ quartet fractions quoted in the
main text refer to the local inelastic response, excluding transitions
within the ground multiplet.

The low-energy effective spin-$1/2$ chain and these local excitations occupy
different scales in the controlled limit.
The effective two-spinon continuum has edges
$\omega_{\mathrm L}(Q)=\pi\Jeff|\sin Q|/2$ and
$\omega_{\mathrm U}(Q)=\pi\Jeff|\sin(Q/2)|$
\cite{descloizeaux1962,caux2006fourspinon}.
Equating its $Q=\pi$ upper edge to the lowest local quartet gap gives formal
values $\beta=0.104,0.176,0.265$ at the three $\alpha$ values.
Their ratios to the cell gap are already $g=1.11,1.14,1.18$.
Such an equality lies outside a small-$g$ expansion and cannot serve as a
controlled prediction of hybridization or of a spectral phase boundary.

\begin{figure*}[t]
 \centering
 \includegraphics[width=0.92\textwidth]{fig5_selection_rules.png}
 \caption{Local projection and selection rules for positive $\alpha$.
 The endpoint coefficients determine the projected exchange, while
 multiplet-resolved weights distinguish bright spin-$3/2$ channels from the
 dark odd-parity doublet at $q=\pi$.
 The connected points in the momentum-weight panel are the five evaluated
 momenta, not a continuous-momentum calculation. No many-body continuum
 boundary or high-energy band assignment is plotted.
 }
 \label{fig:selection}
\end{figure*}

\section{Positive-coupling entropy calibration}
\label{sec:positive-regime}

Entropy offers an independent check of the low-energy interpretation, but the
finite-size slope must be distinguished from a bulk central charge.
At $(\alpha,\beta)=(0.23,0.05)$ we compare the entropy at inter-cell cuts of
$N=8,12,16,20,24$ pentamers ($L=40$--120) with a uniform spin-$1/2$
Heisenberg chain of $N$ effective sites.
The matching is in the number of effective spins, not the physical-site
length. Ratios of the fitted logarithmic slopes are $1.0331$ for the complete
size sequence and $1.0240$ for $N\ge12$; the matched $N=24$ profile gives
$1.0275$ (Fig.~\ref{fig:entanglement}).
The maximum relative entropy change between the tested MPS budgets is
$2.40\times10^{-4}$.
These calibrated values support consistency with the $c=1$ reference.
The raw absolute logarithmic fits fail their fit-quality test and are not
used as precise thermodynamic estimates.

\begin{figure*}[t]
 \centering
 \includegraphics[width=0.94\textwidth]{fig4_entanglement.png}
 \caption{Local-cut entanglement spectra and effective-chain calibration.
 The upper two rows show the lowest relative entanglement levels at five
 inequivalent cuts for $\alpha=0.14$ and $0.33$, both at $\beta=0.5$ and $L=80$.
 The lower comparison is a separate data set at
 $(\alpha,\beta)=(0.23,0.05)$: slopes for pentamer inter-cell cuts are divided
 by those of a uniform spin-$1/2$ chain with the same number of effective
 sites. The three estimators agree with the $c=1$ reference within $3.31\%$.
 The cut spectra are not used as a stand-alone phase classification.
 }
 \label{fig:entanglement}
\end{figure*}

\section{Entropy sensitivity to the local eigensolver}
\label{sec:krylov}
\label{sec:saved-state}

For a general open-chain profile, we fit
\begin{equation}
 S(b)=a_\nu+\frac{\ceff}{6}
 \ln\left[\frac{2L}{\pi}\sin\left(\frac{\pi b}{L}\right)\right],
 \label{eq:entropy}
\end{equation}
with bond-class intercepts $\nu=b\bmod10$ and central windows
$fL\le b\le(1-f)L$, $f=0.10,0.15,0.20,0.25$
\cite{calabrese2004}.
The slope $\ceff$ is descriptive unless finite-size and solver controls
support a central-charge interpretation.

The historical $k=3$ sequence at $(\alpha,\beta)=(-0.10,0.25)$ gives
mean coefficients $0.988,0.695,0.533$ at $L=40,60,80$ for $f=0.20$.
Here $k$ is the local eigensolver's Krylov-subspace dimension, distinct from
the MPS bond dimension $\chi$.
Changing $k=3$ to $8$ at $L=80$ raises the last coefficient to $0.93246$;
the two initial states then differ by only $3.39\times10^{-8}$.
The neighboring $\alpha=-0.175$ control changes from $0.91657$ to $0.95990$.
Figure~\ref{fig:negative-size} keeps these separate fixed-size controls
distinct from the historical size sequence.

At $\alpha=-0.10$, a further $k=8\to12$ control uses both initializations,
64 sweeps, $\chi_{\max}=1200$, and cutoff $10^{-11}$.
The variance per site is $1.9\times10^{-11}$ and the sector gap is
$7.2772\times10^{-5}J_1$.
Across 172 matched observables the changes are bounded by
$7.5\times10^{-12}$ for the gap, $8.4\times10^{-8}$ for the dominant ratio,
and $5.7\times10^{-8}$ for the entropy; $\qpeak=\pi$ is unchanged.
The saved run records and comparison tables have been recovered from the
project archive and checked against their historical checksums.
The anomalously small entropy coefficient does not survive this solver test.
Because the smaller historical sizes still use $k=3$, we do not join them
to the new $L=80$ endpoint as a controlled finite-size extrapolation.

\begin{figure*}[t]
 \centering
 \includegraphics[width=0.96\textwidth]{fig9_entropy_solver_control.pdf}
 \caption{Entropy-coefficient sensitivity at $\beta=0.25$.
 All fits use period-10 bond-class intercepts.
 (a) Historical $k=3$ results at $L=40,60,80$.
 Symbols average the two initial states at $f=0.20$; ranges cover both
 initializations and all four windows.
 (b,c) Separate $L=80$ comparisons of $k=3$ and $8$.
 Symbols show initial-state means at each window and bars the corresponding
 ranges; the two solver series are shifted horizontally by $\mp0.002$ for
 visibility. The dotted line is a reference value, not a bulk central-charge
 determination. Ranges describe deterministic sensitivity, not confidence
 intervals. No $k=12$ values are plotted. The displayed values are taken from the numerical fit tables.}
 \label{fig:negative-size}
\end{figure*}

\section{Spectral resolution and numerical controls}
\label{sec:spectral-control}

The ground-state energies used for the two production spectra are
$E_0/J_1=-12.4925790712$ and $-11.4821383028$, with variance per site
$5.55\times10^{-12}$ and $5.36\times10^{-12}$.
The source norm is checked against the independently contracted equal-time
connected correlations. Complete $L=10$ exact-diagonalization spectra,
processed with the same finite-time estimator, agree with MPS spectra to
relative $L_2$ errors $7.46\times10^{-5}$ and $3.59\times10^{-5}$.
This is a direct small-system validation of the time-evolution and Fourier
normalization, separate from the low-band projection benchmark.

We perform single-variable checks at
$q/\pi=0.2,0.6,1.0$ for both signs of $\alpha$, using the same saved ground
states and $T=60/J_1$.
First, $\chi$ is increased from 192 to 256 at fixed $\Delta t=0.05/J_1$;
then $\Delta t$ is halved at fixed $\chi=256$.
The latter comparison reuses the matching $q=\pi$ record from an earlier
simultaneous refinement of both parameters.
Table~\ref{tab:independent-controls} reports the maximum relative
$L_2$ changes over the three momenta, with the refined spectrum as denominator.
The low-frequency norm is restricted to $0\le\omega\le0.6J_1$.
All 12 comparisons satisfy the criteria fixed before these calculations:
full- and low-frequency changes below $3\%$ and $5\%$, main-peak shifts
below $0.015J_1$, and relative norm drift below $1\%$.
For each of the bands $[0,0.6]$, $[0.6,2.5]$, and $[2.5,4]$ in units of
$J_1$, the weight criterion is
$|W-W_{\rm ref}|\le\max(0.05|W_{\rm ref}|,0.001C_q(0))$.
The largest observed relative band-weight change is $0.376\%$, the largest
norm drift is $0.338\%$, and none of the sampled main peaks changes its
frequency-grid index. These tests constrain the numerical sensitivity of
selected finite-$L$, finite-$T$ spectra; they do not establish long-time,
all-momentum, or thermodynamic convergence.

\begin{table}[!htbp]
 \caption{Independent numerical controls at $L=40$ and $T=60/J_1$.
 Entries are maximum relative $L_2$ changes in percent over
 $q/\pi=0.2,0.6,1.0$. The two columns use $[0,4]$ and $[0,0.6]$ frequency
 intervals in units of $J_1$.}
 \label{tab:independent-controls}
 \begin{ruledtabular}
 \begin{tabular}{rlrr}
 $\alpha$ & Changed parameter & Full & Low\\
 $+0.5$ & $\chi:192\to256$ & 0.0124 & 0.0126\\
 $-0.5$ & $\chi:192\to256$ & 0.0407 & 0.0149\\
 $+0.5$ & $J_1\Delta t:0.05\to0.025$ & 0.8100 & 0.8156\\
 $-0.5$ & $J_1\Delta t:0.05\to0.025$ & 0.5280 & 0.5346\\
 \end{tabular}
 \end{ruledtabular}
\end{table}

The separately computed $L=40$ sector gaps use joint DMRG budgets
$(\chi,\epsilon_{\mathrm{cut}})=(512,10^{-12})$ and $(768,10^{-13})$.
Their changes are of order $10^{-12}J_1$, much smaller than the finite-time
spectral width. Agreement of these two budgets does not make the gaps
thermodynamic quantities or provide a rigorous error estimate.
Similarly, the display step $\Delta\omega=0.005J_1$ does not reduce the
$0.09277J_1$ window width; the heatmaps do not establish sub-linewidth gaps.

For the two heatmaps a single global intensity map is used:
$u(s)=\operatorname{asinh}(s/0.02)/\operatorname{asinh}(6/0.02)$ for
$s=J_1S_T^{zz}\in[0,6]$.
The same map is applied at all $q$ and $\omega$, without energy-dependent
enhancement or rescaling of individual columns.
This is a disclosed color normalization; the accompanying arrays retain
the original linear intensities.
It enhances finite-window sidelobes as well as genuine weak spectral weight,
which is why band identification cannot rest on visual contrast alone.

\subsection{How much weight lies outside the displayed frequency range?}

The finite-time estimator spreads positive-energy lines across zero frequency.
To quantify this effect without modifying the maps, we transform the saved
$C_q(t)$ on a zero-padded discrete Fourier grid with $N_\omega=65536$.
The complete discrete frequency period is
$[-\pi/\Delta t,\pi/\Delta t)$.
The identity $\sum_k S_T^{zz}(q,\omega_k)\Delta\omega=C_q(0)$ is an
algebraic check of this transform, not an independent dynamical sum-rule test
or a proof of infinite-frequency convergence.
An independent contraction of the static connected correlations agrees with
$C_q(0)$ within $5.60\times10^{-11}$ in relative error for nonzero sources.
Direct quadrature on the original display grid reconstructs the saved
spectra with relative error of order $10^{-15}$.

Table~\ref{tab:window-weight} splits the signed spectral integral, normalized
by $C_q(0)$, into negative frequencies, the displayed range, and positive
frequencies above $4J_1$ within the discrete period.
Half the zero-frequency node enters each side; the negative Nyquist node
belongs to the negative side. No negative spectral values are clipped.
Changing $N_\omega$ from 32768 to 65536 changes these fractions by at most
$1.09\times10^{-4}$ over the nonzero-source momentum grid.
This tests numerical quadrature of the same time record, not time resolution.
At $\alpha=-0.5$, the $0.20\pi$ and $0.60\pi$ responses have approximately
$28\%$ and $23\%$ of their weight below zero after window broadening.
These are not physical negative-energy excitations; they demonstrate why
the visible positive-frequency area cannot be identified with the complete
static weight or repaired by stronger color enhancement.

\begin{table}[!htbp]
 \caption{Fractions of the finite-window spectrum, as percentages of $C_q(0)$.
 The intervals refer to the discrete frequency period, with zero split equally
 between the first two columns of weights.}
 \label{tab:window-weight}
 \begin{ruledtabular}
 \begin{tabular}{rrrrr}
 $\alpha$ & $q/\pi$ & $f_-$ & $f_{[0,4]}$ & $f_{>4}$\\
 $+0.5$ & 0.2 & 1.158 & 98.684 & 0.158\\
 $+0.5$ & 0.6 & 0.722 & 99.085 & 0.193\\
 $+0.5$ & 1.0 & 5.683 & 94.168 & 0.149\\
 $-0.5$ & 0.2 & 27.833 & 72.027 & 0.140\\
 $-0.5$ & 0.6 & 22.807 & 77.041 & 0.152\\
 $-0.5$ & 1.0 & 6.798 & 93.022 & 0.180\\
 \end{tabular}
 \end{ruledtabular}
\end{table}
